\documentclass{article}

\usepackage[utf8]{inputenc}
\usepackage{amsmath, esint}
\usepackage{wasysym}
\usepackage{qrcode}
\usepackage[colorlinks]{hyperref}
\usepackage{lmodern}
\usepackage{graphicx}
\usepackage{xcolor}
\usepackage[left=2cm, top=3cm, right=2cm]{geometry}
\usepackage{minted}
\usepackage{booktabs}
\usepackage{svg}
\usepackage{xcolor}
\definecolor{LightGray}{gray}{0.975}
\hypersetup{
  colorlinks,
  linkcolor=blue,
  urlcolor=blue
}

\title{ML101 \\ Homework 02: A `gentle' Introduction to R.}
\author{Andrés Oswaldo Calderón Romero, Ph.D.}
\date{\today}

\begin{document}

\maketitle

\section{Introduction}
\label{sec:intro}

Welcome to Homework 02. This assignment serves as a practical, gentle introduction to R and its computing environment. Before diving into advanced statistical modeling and machine learning algorithms, it is essential to build familiarity with fundamental commands, data structures, and workspace setups. This lab is designed to get your local environment configured correctly and to provide hands-on practice with the basic tools you will use throughout the semester.

\section{Setting up}
\label{sec:setup}

Follow this brief tutorial to set up an R environment on your preferred operating system. The guide is adapted from Garrett Grolemund’s book, \href{https://rstudio-education.github.io/hopr/}{\textit{Hands-On Programming with R: Write Your Own Functions and Simulations}}. You can access the \href{https://rstudio-education.github.io/hopr/starting.html}{setup tutorial here}.

\section{A Brief R Introduction}
\label{sec:rinto}

We will now complete the lab exercise from Chapter 2 (pages 42–52) of \href{https://www.statlearning.com/}{\textit{An Introduction to Statistical Learning: with Applications in R}} by Gareth James, Daniela Witten, Trevor Hastie, and Robert Tibshirani. This lab will introduce essential R commands and ensure your environment is properly set up for upcoming hands-on exercises.

\section{Individual Assignment}
\label{sec:work}

No submission is required for this exercise; however, you should ensure your R environment is working properly so you are prepared to use these essential commands in future labs.

\section*{Optional: A Step into Linear Regression}
To prepare for our next topic, we will use an excellent resource provided by Stanford Online. We will watch a lecture by Prof.~Christopher Ré from the CS229 Machine Learning course (taught in Spring 2026). For our purposes, please watch from minute 5:51 to 32:33 of the \href{https://youtu.be/cmNIMjPYdgM?t=351}{CS229 Lecture 2 video}\footnote{Since this is a YouTube video, remember that you can activate Spanish subtitles if you prefer.}.

\vspace{5mm}
Happy Hacking! \includesvg[width=4mm]{figures/sunglasses}

\end{document}

